\color{uniblue}\section[Выбор уставок защиты от перегрузки трансформатора]{\sloppy Выбор уставок защиты от перегрузки \mbox{трансформатора}} 
\color{black}

\begin{enumerate}[label=\arabic{section}.\arabic{enumi}, labelsep=7pt, leftmargin=0pt, itemindent=60pt, align=left, nosep]

\item
Ток срабатывания защиты от перегрузки стороны высшего (низшего) напряжения трансформатора определяется по выражению:
\begin{equation}
    \label{eq:zp1}
    I_{\textsl{с.з.ЗП}} = \frac{K_{\textsl{отс}}}{K_{\textsl{в}}} \cdot I_{\textsl{ном,ст}},
    \end{equation}
    \begin{eqexpl}[15mm]
    \item{$I_{\textsl{ном.ст.}}$} номинальное значение тока стороны трансформатора, А;
    \item{$K_{\textsl{отс}}$} коэффициент отстройки, принимается равным 1,05;
    \item{$K_{\textsl{в}}$} коэффициент возврата, принимается равным 0,95.
    \end{eqexpl}

\item
Защита от перегрузки при наличии обслуживающего персонала на подстанции выполняется с действием на сигнал, выдержку времени в этом случае выбирают из условия отстройки от времени самозапуска электродвигателей, подключенных к секциям шин 6-10 кВ ($t_{\textsl{ср}} \ge 10 \textsl{с}$). 
\item
При отсутствии обслуживающего персонала на подстанции возможно выполнение защиты от перегрузки с действием на отключение трансформатора, в этом случае выдержка времени на отключение трансформатора выбирается по перегрузочной характеристике трансформатора в соответствии с \cite{bib:overload}.

\end{enumerate}


